\begin{abhakseite}
\abhakgruppe{Einheit 1 · Monotonie und erste Ableitung}
\abhakgruppe{\normalfont\itshape Monotonie aus der Ableitung bestimmen}
\abhak{9}{Ich kann am Vorzeichen der Ableitung ablesen, ob der Graph steigt oder fällt.}
\abhak{10}{Ich kann die Monotonieintervalle einer Funktion bestimmen.}
\abhakgruppe{\normalfont\itshape Monotonie am Term nachweisen}
\abhak{11}{Ich kann am Term der Ableitung nachweisen, dass ein Graph überall steigt.}
\abhak{12}{Ich finde den Fehler bei Monotonieintervallen.}
\abhak{13}{Ich kann begründen, ob ein Graph überall steigt.}
\abhak{14}{Ich kann die Monotoniebereiche zweier Modelle vergleichen.}
\abhakgruppe{Einheit 2 · Extrempunkte}
\abhak{15}{Ich erkenne, ob eine Stelle gegeben ist oder gesucht wird.}
\abhak{16}{Ich erkenne, welcher Nachweis für die Art des Punktes passt.}
\abhakgruppe{\normalfont\itshape Extrempunkte berechnen}
\abhak{17}{Ich kann Extrempunkte berechnen.}
\abhak{18}{Ich kann Extrempunkte berechnen – weiter.}
\abhakgruppe{\normalfont\itshape Extrempunkte nachweisen}
\abhak{19}{Ich kann nachweisen, dass an einer gegebenen Stelle ein Extrempunkt liegt.}
\abhak{20}{Ich finde den Fehler bei der Bestimmung von Extrempunkten.}
\abhak{21}{Ich kann Aussagen über Extrempunkte begründen.}
\abhak{22}{Ich kann mit Extrempunkten eine Bedingung im Sachzusammenhang prüfen.}
\abhakgruppe{Einheit 3 · Krümmung und Wendepunkte}
\abhakgruppe{\normalfont\itshape Wendepunkte berechnen}
\abhak{23}{Ich kann Wendepunkte berechnen.}
\abhak{24}{Ich kann angeben, wo ein Graph links- und wo er rechtsgekrümmt ist.}
\abhakgruppe{\normalfont\itshape Wendepunkte nachweisen}
\abhak{25}{Ich kann nachweisen, dass an einer gegebenen Stelle ein Wendepunkt liegt.}
\abhak{26}{Ich finde den Fehler bei Krümmung und Wendepunkten.}
\abhak{27}{Ich kann begründen, wann ein Wendepunkt vorliegt.}
\abhak{28}{Ich kann Krümmung im Sachzusammenhang deuten.}
\abhakgruppe{Einheit 4 · Graph und Ableitungsgraph}
\abhak{29}{Ich kann am Graphen von $f$ ablesen, wo $f'$ positiv, negativ oder null ist.}
\abhak{30}{Ich kann am Graphen von $f'$ ablesen, was der Graph von $f$ macht.}
\abhak{31}{Ich kann die Graphen von $f$ und $f'$ aus berechneten Werten zeichnen.}
\abhak{32}{Ich kann einen möglichen Graphen zu gegebenen Eigenschaften skizzieren.}
\abhak{33}{Ich kann Aussagen über Graphen und Ableitungen beurteilen.}
\abhak{34}{Ich finde den Fehler beim Lesen von Graph und Ableitungsgraph.}
\abhak{35}{Ich kann Zusammenhänge zwischen $f$ und $f'$ begründen.}
\abhak{36}{Ich kann aus dem Graphen einer Änderungsrate auf den Bestand schließen.}
\abhakgruppe{Einheit 5 · Kurvenuntersuchung im Sachzusammenhang}
\abhak{37}{Ich erkenne, was eine Sachfrage mathematisch bedeutet.}
\abhakgruppe{\normalfont\itshape Verlauf am Graphen beschreiben}
\abhak{38}{Ich kann den Verlauf eines Graphen im Sachzusammenhang beschreiben.}
\abhakgruppe{\normalfont\itshape Größter und kleinster Wert}
\abhak{39}{Ich kann den größten Wert im Sachzusammenhang berechnen.}
\abhakgruppe{\normalfont\itshape Stärkste Zunahme und Abnahme}
\abhak{40}{Ich kann berechnen, wann eine Größe am stärksten zu- oder abnimmt.}
\abhakgruppe{\normalfont\itshape Maße aus einem Profil}
\abhak{41}{Ich kann Maße eines Profils mit einem Maßstab berechnen.}
\abhak{42}{Ich finde den Fehler bei Sachaufgaben zur Kurvenuntersuchung.}
\abhak{43}{Ich kann Ergebnisse im Sachzusammenhang begründen.}
\abhak{44}{Ich kann den Wendepunkt im Sachzusammenhang deuten.}
\end{abhakseite}
